\documentclass[10pt,a4paper]{amsart}
\usepackage[margin=19mm]{geometry}
\usepackage{amsmath,amssymb,mathtools}
\usepackage[scr=rsfs,cal=euler]{mathalfa}
\usepackage{xfrac}
\usepackage[unicode,hidelinks,bookmarks=false]{hyperref}
\newcommand{\ie}{i.e.\ }
\newcommand{\cf}{cf.\ }
\newcommand{\resp}{respectively\ }
\newcommand{\Subsection}{\subsection}
\providecommand{\nobreakdash}{\nobreak\hbox{-}}
\setlength{\emergencystretch}{2em}
\multlinegap0pt
\newtheorem{comment}{Comment}
\usepackage[shortlabels]{enumitem}
\usepackage{bm}
\usepackage{amssymb}
\usepackage{algorithm}\usepackage{algpseudocode}
\newtheorem{thm}{Theorem}
\newcommand{\R}{\mathbb{R}}
\newcommand{\bff}{\boldsymbol{f}}
\newcommand{\bx}{\boldsymbol{x}}
\title{Triply-Scalable Equivariant Gaussian Process Modeling}
\author{Tim Steinert, David Ginsbourger}
\date{Source: arXiv:2609.21085v1 (CC BY 4.0)}
\begin{document}
\maketitle

\noindent\emph{Source and licence.} Triply-Scalable Equivariant Gaussian Process Modeling. Version arXiv:2609.21085v1 (CC BY 4.0), distributed by arXiv under the Creative Commons Attribution 4.0 International licence, http://creativecommons.org/licenses/by/4.0/, as recorded in the arXiv licence metadata for this exact version and shown on the abstract page https://arxiv.org/abs/2609.21085v1. Full licence notice: \texttt{licenses/arxiv-2609.21085-CC-BY-4.0.txt}.\par\smallskip

\noindent\textit{Editorial scope.} Counted unit: the article's thm thm:equivariant-conditioning (source0.tex lines 489--501). The article's complete original proof of it is delivered with it (source0.tex lines 503--555, one \texttt{proof} environment of 53 lines). No mathematics is written, repaired or re-derived: the statement and the proof are the article's own lines, in the article's own notation, and no retained line is substituted or re-flowed.  No further source-local context is needed: every cross-reference of every retained line resolves inside the two delivered spans.  Nothing was omitted from any retained span, so the package carries no omission record.  No retained line contains a citation, so the wrapper carries no bibliography and no bibliography entry.  No retained line contains a figure environment, an image-inclusion command or drawing code, so no image is delivered and the package declares no asset dependency.  Editorial formatting only, in addition: an AMS \texttt{amsart} preamble that restates the article's own declarations and macros used by the retained lines, namely the environments \texttt{thm} on separate counters, together with the standard packages those lines need.  The retained environments print in this document as Theorem 1 in order of appearance, whereas the article prints the same items with the numbering of its own full document.  The complete native source of the article as distributed in the arXiv e-print for this exact version is bundled unchanged as \texttt{sources/210934/source0.tex}.
\begin{thm}[Equivalences between unconditional and conditional stochastic equivariance of GPs]\label{thm:equivariant-conditioning}
Let $\bff$ be a centred $\mathbb{R}^{p}$-valued Gaussian process on $(\Omega, \mathcal{A}, \mathbb{P})$ indexed by $D$ with kernel $K \colon D \times D \to \R^{p \times p}$. Let a linear group $G$ act on $D$ via $\star$ and on $\R^{p}$ by a representation $\rho \colon g \mapsto \rho_{g} \in \R^{p \times p}$.
Then the following are equivalent:
\begin{enumerate}[label=(\roman*)]
    \item For all
    $\bx \in D$ and $g \in G$,
    \begin{equation*}\mathbb{P}\big(\bff(g \star \bx) = \rho_{g} \bff(\bx)\big) = 1.\end{equation*}
    \item For \textbf{any} finite collection of inputs in $D,$ $X = \left(\bx_{1},\dots,\bx_{n}\right)$ and $\R^{pn}$-valued random vector $\boldsymbol{\varepsilon}$ on $(\Omega, \mathcal{A}, \mathbb{P}),$ for all $\bx \in D$ and $g \in G$,
    \begin{equation*}\mathbb{P}\big(\bff(g \star \bx) = \rho_{g} \bff(\bx) \,\big|\, \bff(X)+\boldsymbol{\varepsilon}\big)= 1 \quad \text{almost surely.}\end{equation*}
    \item For \textbf{some} finite collection of inputs in $D,$ $X = \left(\bx_{1},\dots,\bx_{n}\right),$ and $\R^{pn}$-valued random vector $\boldsymbol{\varepsilon}$ on $(\Omega, \mathcal{A}, \mathbb{P}),$ for all $\bx \in D$ and $g \in G$,
    \begin{equation*}\mathbb{P}\big(\bff(g \star \bx) = \rho_{g} \bff(\bx) \,\big|\, \bff(X)+\boldsymbol{\varepsilon}\big)= 1 \quad \text{almost surely.}\end{equation*}
\end{enumerate}
\end{thm}

\begin{proof}
\begin{comment}
    $(i)\Rightarrow (ii):$
    For $x\in D,$ and any a finite subset $X \subset D,$ define the random vectors $\boldsymbol{A}:=\bff(g\star x)-\rho_g\bff(x)$ and $\boldsymbol{B}:=\bff(X).$ Since $\bff$ is a Gaussian process on $D$ and $\boldsymbol{A}$ is a linear transformation of $\bff,$ $(\boldsymbol{A},\boldsymbol{B})$ is jointly Gaussian. It is thus sufficient to show that \begin{equation*}\mathbb{E}_{\boldsymbol{B}}\left[\mathbb{E}_{\boldsymbol{A}\mid \boldsymbol{B}}\left[\|\boldsymbol{A}\|^2\right]\right]=0,\end{equation*}
    as $0\leq \mathbb{E}_{\boldsymbol{A}\mid \boldsymbol{B}}\left[\|\boldsymbol{A}\|^2\right]$ then implies that $\mathbb{P}(\boldsymbol{A}=0\mid \boldsymbol{B})=1$ almost surely. Now,
    \begin{align*}
        \mathbb{E}_{\boldsymbol{B}}\left[E_{\boldsymbol{A}\mid \boldsymbol{B}}\left[\|\boldsymbol{A}\|^2\right]\right]&=\mathbb{E}_{\boldsymbol{B}}\left[ \mathbb{E}_{\boldsymbol{A}\mid \boldsymbol{B}}\left[\operatorname{tr}\left(\boldsymbol{A}^\top \boldsymbol{A}\right)\right]\right]\\
        &=\mathbb{E}_{\boldsymbol{B}}\left[ \mathbb{E}_{\boldsymbol{A}\mid \boldsymbol{B}}\left[\operatorname{tr}\left(\boldsymbol{A} \boldsymbol{A}^\top\right)\right]\right]\\
        &=\operatorname{tr}\left(\mathbb{E}_{\boldsymbol{B}}\left[\mathbb{E}_{\boldsymbol{A}\mid \boldsymbol{B}}\left[\boldsymbol{A} \boldsymbol{A}^\top\right]\right]\right).
    \end{align*}
    Since
    \begin{align*}
        \mathbb{E}_{\boldsymbol{A}\mid \boldsymbol{B}=\boldsymbol{b}}\left[\boldsymbol{A} \boldsymbol{A}^\top\right]&=\operatorname{Cov}(\boldsymbol{A}\mid \boldsymbol{B}=\boldsymbol{b})+ \mathbb{E}[\boldsymbol{A}\mid \boldsymbol{B}=\boldsymbol{b}]\mathbb{E}[\boldsymbol{A}\mid \boldsymbol{B}=\boldsymbol{b}]^\top\\
        &=K_{\boldsymbol{A}}-K_{\boldsymbol{A}\boldsymbol{B}}K_{\boldsymbol{B}}^{-1}K_{\boldsymbol{B}\boldsymbol{A}}+ K_{\boldsymbol{A}\boldsymbol{B}}K_{\boldsymbol{B}}^{-1}\boldsymbol{b}\boldsymbol{b}^\top K_{\boldsymbol{B}}^{-1}K_{\boldsymbol{B}\boldsymbol{A}},
    \end{align*}
    where $K_{\boldsymbol{A}\boldsymbol{B}}$ denotes $\operatorname{Cov}(\boldsymbol{A},\boldsymbol{B}),$ we obtain by the zero mean assumption
    \begin{align*}
        \mathbb{E}_{\boldsymbol{B}}\left[E_{\boldsymbol{A}\mid \boldsymbol{B}}\left[\boldsymbol{A}\boldsymbol{A}^\top\right]\right]&=\mathbb{E}_{\boldsymbol{B}}\left[K_{\boldsymbol{A}}-K_{\boldsymbol{A}\boldsymbol{B}}K^{-1}_{\boldsymbol{B}}K_{\boldsymbol{B}\boldsymbol{A}}+ K_{\boldsymbol{A}\boldsymbol{B}}K^{-1}_{\boldsymbol{B}}\boldsymbol{b}\boldsymbol{b}^\top K_{\boldsymbol{B}}^{-1}K_{\boldsymbol{B}\boldsymbol{A}}\right]\\
        &=K_{\boldsymbol{A}}-K_{\boldsymbol{A}\boldsymbol{B}}K^{-1}_{\boldsymbol{B}}K_{\boldsymbol{B}\boldsymbol{A}}+ K_{\boldsymbol{A}\boldsymbol{B}}K^{-1}_{\boldsymbol{B}}K_{\boldsymbol{B}\boldsymbol{A}}\\
        &=K_{\boldsymbol{A}}.
    \end{align*}
    The direction follows as $\mathbb{P}(\boldsymbol{A}=0)=1$ implies $\operatorname{tr}(K_{\boldsymbol{A}})=0.$
\end{comment}

Fix $\bx\in D$ and $g\in G$, and define the event $E:=\{\bff(g\star\bx)-\rho_g\bff(\bx)=\bm 0\}$. By assumption (i), $\mathbb{P}(E)=1$. Now take any finite collection $X = \left(\bx_{1},\dots,\bx_{n}\right)$ in $D$ and any $\R^{pn}$-valued random vector $\boldsymbol{\varepsilon}$ on $(\Omega, \mathcal{A}, \mathbb{P})$.
Then $\mathbb{P}(E\mid \bff(X)+\boldsymbol{\varepsilon})=1$ almost surely since $\mathbb{E}[\mathbb{P}(E\mid \bff(X)+\boldsymbol{\varepsilon})]=\mathbb{P}(E)=1$ by the law of total probability, and thus $\mathbb{P}(E\mid \bff(X)+\boldsymbol{\varepsilon})<1$ with positive probability would result in a contradiction. This implies (ii) and in turn (iii). Now, assuming (iii), we again obtain by the law of total probability that
\begin{equation}
\mathbb{P}(E)=\mathbb{E}\big[\mathbb{P}(E\mid \bff(X)+\boldsymbol{\varepsilon})\big]=1.
\end{equation}
%and equivalently $\mathbb P(A\neq 0)=0.$ Denote the $\sigma-$algebra $\mathcal{G}:=\sigma(f(X))$ which is generated by cylindrical sets of the form
%\begin{equation*}%C_B:=\{\omega:\bff(X)(\omega)\in B\}, %\qquad B\in \mathcal B((\mathbb R^p)^n). %\end{equation*} To prove
%\begin{equation*}% \mathbb P(\boldsymbol{A}=0\mid \mathcal G)=1 \qquad \text{a.s.}, %\end{equation*}
%it is enough to show that
%\begin{equation*}% \mathbb P(A\neq0\mid \mathcal G)=0 \qquad \text{a.s.} %\end{equation*}
%By definition, the random variable
%\begin{equation*}%Z:=\mathbb P(A\neq 0\mid \mathcal G) %=\mathbb E[\mathbbm{1}_{\{A\neq 0\}}\mid \mathcal G] %\end{equation*}
%is $\mathcal G$-measurable and satisfies
%\begin{equation*}%\int_{C_B} Z\,d\mathbb P %= %\mathbb P((A\neq 0)\cap C_B) %\qquad\text{for all }C_B\in\mathcal G. %\end{equation*}
%For every such cylindrical set,
%\begin{equation*}%\mathbb P((A\neq 0)\cap C_B)\le \mathbb P(A\neq 0)=0, %\end{equation*}
%hence
%\begin{equation*}%\int_{C_B} Z\,d\mathbb P=0. %\end{equation*}
%Since the cylindrical sets generate $\mathcal G$, the same identity holds for all $C\in\mathcal G$. In particular, for
%\begin{equation*}%C=\{Z>0\}\in\mathcal G, %\end{equation*}
%we obtain
%\begin{equation*}%\int_{\{Z>0\}} Z\,d\mathbb P=0. %\end{equation*}
%As $Z\ge 0$ almost surely, this implies $Z=0$ almost %surely.
%    $(ii)\Rightarrow(i):$
%Fix $x\in D$ and $g\in G$. For any finite $X\subset D$, by %the tower property,
%\begin{equation*}%\mathbb P(\boldsymbol{A}=0) %=\mathbb E[\mathbf 1_{\{\boldsymbol{A}=0\}}] %=\mathbb E\!\left[\mathbb E[\mathbf 1_{\{\boldsymbol{A}=0\}}\mid %\bff(X)]\right] %=1, %\end{equation*}
%where the last equality follows from (ii). Thus
%\begin{equation*}%\mathbb P\bigl(\bff(g\star x)=\rho_g\bff(x)\bigr)=1. %\end{equation*}
\end{proof}

\end{document}
